\documentclass[../main.tex]{subfiles}

\begin{document}
\section{Background} 
\begin{lemma}
  \label{lem:structure-of-roots}
  Let $f \in \Z\left[ x \right]$, and let $L$ be its splitting field over $\Q$.
  Assume that over $\F_p$, the polynomial $\overline{f}$ decomposes as
  \[
    \overline{f}\left( x \right) = \prod_{i=1}^{k} \overline{f_i}\left( x \right)^{e_i}
  \]
  where the $\overline{f_i}$ are distinct monic irreducible polynomials.
  Let $\P$ be a prime of $\O_L$ lying over $\left( p \right) \subseteq \Z$.
  Then there is a partition, of the set of roots of $f$ into disjoint subsets $S_1, \ldots, S_k$ such that: 
  \begin{enumerate}
    \item for each $1 \leq i \leq k$, $\left| S_i \right| = e_i \deg\left( \overline{f_i} \right)$
    \item The map 
    \[
      \pi_{\P} : \O_L \to \O_L / \P
    \]
    reduces to a surjective $e_i$ to $1$ map from $S_i$ onto the set of roots of $\overline{f_i}$ over $\F_p$.
    \item The decomposition group at $\P$, $D\left( \P \right)$, preserves each $S_i$ 
    \item The decomposition group at $\P$ acts on $S_i$ in a way that respects the partition into fibers of $\pi_{\P}$.
    \item The inertia group at $\P$ acts trivially on the set of fibers of $\pi_{\P}$.
    \item The preimage of the frobenious automorphism $\phi : x \mapsto x^p$ under $\pi_{\P}$ acts on $S_i$ that respects the partition into fibers of $\pi_{\P}$, and induces on the set of fibers a cycle of length $\deg\left( \overline{f_i} \right)$.
  \end{enumerate}
\end{lemma}
\begin{proof}
  By Hensel's lemma there are $f_1, \dots, f_k \in \Z_p\left[ x \right]$ such that for each $1 \leq i \leq k$, $\left( \overline{f_i} \right)^{e_i}$ is the image of $f_i$ under the cannonical map $\Z_p\left[ x \right] \to \F_p\left[ x \right]$, and
  \[
    f\left( x \right) = \prod_{i=1}^{k} f_i\left( x \right)
  \]
  over $\Z_p$.
  Set $S_i$ to be the set of roots of $f_i$ in $L$.

  Fix $1 \leq i \leq k$ and let $\beta_1, \ldots, \beta_{\deg\left( \overline{f_i} \right)}$ be the roots of $\overline{f_i}$ over $\F_p$.
  Then 
  \begin{equation}
    \label{eq:projections_of_roots}
    \begin{aligned}
      & \prod_{\alpha \in S_i} \left( x - \pi_{\P}\left( \alpha \right) \right) = \\
      & \pi_{\P}\left( \prod_{\alpha \in S_i} \left( x - \alpha \right) \right) = \\
      & \pi_{\P}\left( f_i\left( x \right) \right) = \\
      & \overline{f_i}\left( x \right)^{e_i} = \\
      & \prod_{j=1}^{\deg\left( \overline{f_i} \right)} \left( x - \beta_j \right)^{e_i}
    \end{aligned}
  \end{equation}
  and hence, since $\F_p^\alg\left[ x \right]$ is a unique factorization domain, for each $1 \leq j \leq \deg\left( \overline{f_i} \right)$, there are exactly $e_i$ elements $\alpha \in S_i$ such that $\pi_{\P}\left( \alpha \right) = \beta_j$.
  This proves (2) and in particular it implies (1).

  Let $\sigma \in D\left( \P \right)$.
  Then $\sigma$ extends to an automorphism $\widehat{\sigma}$ of the completion of $L$ with respect the $\P$-valuation $v_\P$.
  We shall denote this completion by $\widehat{L}$.
  Since $\widehat{\sigma}$ fixes $\Q$, it also fixes its closure with respec to $v_\P$, a closure which is also known as $\Q_p$.
  In particular, the automorphism it induces on $\widehat{L}\left[ x \right]$ fixes $\Z_p\left[ x \right]$ and thus it fixes $f_i$ and consequently so does the automorphism that $\sigma$ it induces $L\left[ x \right]$.
  Therefore,  
  \begin{align*}
    & \prod_{\alpha \in S_i} \left( x - \sigma\left( \alpha \right) \right) = \\
    & \sigma\left( \prod_{\alpha \in S_i} \left( x - \alpha \right) \right) = \\
    & \sigma\left( f_i \right) = \\
    & f_i = \\
    & \prod_{\alpha \in S_i} \left( x - \alpha \right)
  \end{align*}
  and hence, since $L\left[ x \right]$ is a unique factorization domain, it follows that $\sigma\left( S_i \right) = S_i$.
  Since the same is true for any $i$ and for any $\sigma \in D\left( \P \right)$, it follows that $D\left( \P \right)$ preserves each $S_i$ and.

  It is well known that each element of the decomposition group induces an automorphism of the $L/\P$, which in particular implies that it respects the partition of $L$, and in particular of $S_i$ into fibres of $\pi_\P$.
  The automorphism that any element of the inertia group, $I\left( \P \right)$, induces on $L/\P$ is the identity, by the mere definition of the inertia group, and hence the bijection that any such automorphism induces from the set of fibres to iteself is also the identity.

  By a slight abuse of notation, use $\phi$ to denote some member of $D\left( \P \right)$ that induces the frobenious map $\phi$ on $L/\P$ (it is a well known fact that any automorphism of $L/\P$ is induces by some automorphism of $L$).
  Then as we have already shown, sicne it is a member of $D\left( \P \right)$ its action on $S_i$ preserves the partition into fibres of $\pi_\P$.
  Moreover, the action it induces on them is the action it induces on their images in $L / \P$ which is indeed a cycle of order 
  \[
    \left[ L/\P : \F_p \right] = \deg\left( \overline{f_i} \right)
  \]
\end{proof}
\begin{definition}
  In the situation of lemma \ref{lem:structure-of-roots}, we will call the each $S_i$ a decomposition block of $\P$ and we will call the intersection of each $S_i$ with each fibre under $\pi_\P$ an inertia block of $\P$.
\end{definition}


\end{document}